\label{chap:literature}
This chapter reviews the work on which the thesis builds and against which it is measured. It moves from the
navigation problem itself (\cref{sec:lit-nav,sec:lit-objnav}) to the use of foundation models as embodied agents
(\cref{sec:lit-agents}), and then to three ingredients of \ours: memory that attaches semantics to geometry
(\cref{sec:lit-memory}), verification and search with VLM queries (\cref{sec:lit-verify}), and geometry from images
(\cref{sec:lit-geometry}). \Cref{sec:lit-gap} states the gap that the thesis addresses.

\section{Embodied navigation}
\label{sec:lit-nav}
\subsection{Classical navigation}
Classical mobile robots separate mapping, localization, and planning. Occupancy grids represent free and occupied
space as a probabilistic grid built from range measurements~\cite{elfes1989occupancy}; frontier-based exploration
drives the robot to the boundary between known free space and unknown space until the map is complete~\cite{yamauchi1997frontier};
and graph search such as A*~\cite{hart1968astar} computes collision-free paths on the map. These components are
reliable, interpretable, and require no learning, and they form the geometric core of most modular navigation
systems~\cite{thrun2005probabilistic}. What they lack is semantics: a classical map records where the robot can go,
not where a ``kettle'' or ``something to sit on'' is.

\subsection{Benchmarks for embodied navigation}
Simulators of indoor scenes, such as AI2-THOR~\cite{kolve2017ai2thor} and Habitat~\cite{savva2019habitat}, made
large-scale evaluation of embodied agents possible. Anderson~\etal{} proposed the evaluation framework that most
benchmarks follow, including success weighted by path length (SPL)~\cite{anderson2018evaluation}. Object-goal
navigation (ObjectNav) asks the agent to reach an instance of a named category~\cite{batra2020objectnav};
RoboTHOR~\cite{deitke2020robothor} poses it in apartments built both in simulation and in the real world.
Vision-and-language navigation (VLN) asks the agent to follow route instructions~\cite{anderson2018vln}.
EmbodiedBench~\cite{yang2025embodiedbench} evaluates multimodal LLMs on high- and low-level embodied tasks through a
standardized action interface, and its EB-Navigation benchmark is the main testbed of this thesis. Related benchmarks
probe embodied decision making and navigation reasoning of language models~\cite{li2024embodied,qiao2025navbench}.

\subsection{Learned navigation policies}
End-to-end learning maps observations directly to actions. Policies trained for ObjectNav in large procedurally
generated houses reach strong results on RoboTHOR~\cite{deitke2022procthor}. Semantic exploration learns where a
category is likely to be and builds an explicit semantic map~\cite{chaplot2020object}. More recently,
vision-language-action models fine-tune a VLM to output navigation actions from video~\cite{zhang2024navid,cheng2024navila,wei2025streamvln},
and diffusion policies and world models generate trajectories or predict future views for planning~\cite{sridhar2023nomad,cai2025navdp,bar2024navigation}.
Learned policies capture regularities of indoor scenes that hand-written rules miss, such as where kitchens tend
to be. They pay for this with task-specific data, long training on large simulated datasets, and sensitivity to the
gap between the training scenes and the deployment scenes. \ours uses a frozen model without training and keeps the
navigation itself in classical geometry, which transfers without retraining.

\subsection{Compact vision-language models}
The VLMs that a robot can run locally are the small members of open model families. Qwen2.5-VL and Qwen3-VL provide
models of 3--8B parameters that answer questions about images and return bounding boxes for objects named in
text~\cite{bai2025qwen25vl,bai2025qwen3vl}. Weight quantization such as AWQ stores the weights in 4 bits with a small
loss in accuracy~\cite{lin2024awq}, and serving engines with paged attention run such models efficiently on a single
GPU~\cite{kwon2023vllm}. Together they make it possible to run a 4B model next to a simulator on one 8\,GB consumer
GPU, the setting of this thesis. On embodied benchmarks, however, these models lag far behind frontier models when
they are used as agents~\cite{yang2025embodiedbench}; the thesis asks whether this gap reflects their perception or
their role.

\section{Zero-shot object navigation}
\label{sec:lit-objnav}
Zero-shot ObjectNav removes task-specific training by combining an open-vocabulary perception model with a map and a
planner. CoW pairs open-vocabulary localization with frontier exploration~\cite{gadre2022cows}. Later systems add
semantic reasoning about where to search. ESC and L3MVN ask a language model which frontier is likely to lead to
the target~\cite{zhou2023esc,yu2023l3mvn}; VLFM builds a value map from vision-language similarity~\cite{yokoyama2023vlfm};
OpenFMNav and LGX use foundation models to propose and select goals~\cite{kuang2024openfmnav,dorbala2024lgx};
SG-Nav, MSGNav, and UniGoal reason over scene or goal graphs with an LLM~\cite{yin2024online,huang2025msgnav,yin2025unigoal};
WMNav uses a world model~\cite{nie2025wmnav}; CogNav models navigation as a cognitive state machine~\cite{cao2024cognav};
STRIVE queries a VLM over a structured representation~\cite{zhu2025strive}; and MVP-Nav and ForesightNav plan over
multi-layer value maps and semantic-geometric predictions~\cite{xie2026mvpnav,liu2026foresightnav}.
Instruction-following agents extend this structure from categories to free-form instructions by grounding
language in maps and scene graphs~\cite{long2024instructnav,rana2023sayplan,werby2024hierarchical,huang2023vlmaps,shah2022lmnav}.

These systems share the modular structure of \ours: perception produces detections, a map stores them, and a
planner acts. They differ from \ours in three respects. First, most use dedicated open-vocabulary detectors such as
OWL-ViT or Grounding DINO~\cite{minderer2022owlvit,liu2024groundingdino}, often together with a large language model
for reasoning; \ours uses one small generative VLM as its only perception model, which reads open instructions
directly. Second, many ask a language model where to explore; \ours never asks the model where to go, and its
exploration follows geometric rules. Third, a detector returns calibrated scores for a fixed vocabulary, whereas a
generative VLM writes its boxes as text and places boxes in views that do not show the target
(\cref{sec:diag-probe}). The memory of \ours is designed for this error profile.

\section{Multimodal LLMs as embodied agents}
\label{sec:lit-agents}
A second line of work uses a multimodal LLM as the whole agent. SayCan grounds an LLM's plans in the affordances
of learned skills~\cite{ahn2022saycan}. In EmbodiedBench, the model maps the image, the instruction, and the text
history directly to a plan of actions~\cite{yang2025embodiedbench}. Agents are improved by embodied pre-training,
online reinforcement learning, or policy-level reasoning~\cite{chen2025transforming,wang2025vagen,zhang2025embodiedx,liu2025robogpt},
or at test time with reflection, experience, workflow memory, multi-memory architectures, and
verifiers~\cite{shinn2023reflexion,zhao2023expel,chen2024automanual,wang2024agent,lei2025robomemory,singhi2026think}.

On EB-Navigation specifically, training-free agents add a scene graph~\cite{huang2025esca}, a module graph~\cite{wang2026mosaic},
a skill library~\cite{ju2026embodiskill}, brain-inspired memory~\cite{ma2026brainmem},
an optimized harness~\cite{lee2026emho}, self-evolution~\cite{liang2026robofoundry}, or a library of perception and
mapping tools~\cite{zhou2026enabling}, and SeeNav-Agent renders the target into a bird's-eye view~\cite{wang2025seenav}.
Fine-tuned agents train world models or a Q-model around a 4B policy~\cite{jang2026musix,xie2026sva}.
All of these agents keep the model as the planner and as the reader of its own memory. Whatever is added, a
scene graph, a skill, or a tool output, enters the model's context as text or as an image, and the model must
interpret it and decide. This places the hardest part of the task, integrating observations over time into a
decision, on the model. The reported gains therefore grow with model size: the training-free agents on
EB-Navigation use models of 8B parameters up to frontier scale, and the agents that work with 4B models require
training. \ours takes the opposite position: the VLM is only a sensor, and its answers feed an explicit metric
estimator and planner.

\Cref{tab:lit-compare} contrasts the main families of approaches along the dimensions that matter for this thesis.

\begin{table}[htbp]
\centering\small
\caption[Families of approaches to instruction-following navigation]{Families of approaches to instruction-following
navigation. ``Open'' instructions describe the target by use or appearance as well as by name.}
\label{tab:lit-compare}
\resizebox{\linewidth}{!}{%
\begin{tabular}{@{}lllllc@{}}
\toprule
Family & Who decides & Perception & Spatial state & Training & Open instructions\\
\midrule
Prompted MLLM agent~\cite{yang2025embodiedbench} & the model & the model & text history & none & yes\\
Training-free MLLM agents~\cite{lee2026emho,ma2026brainmem,zhou2026enabling} & the model & model, tools & text, graphs, tool state & none & yes\\
Fine-tuned VLM agents~\cite{xie2026sva,jang2026musix} & the model & the model & learned & task-specific & yes\\
Modular zero-shot ObjectNav~\cite{gadre2022cows,yokoyama2023vlfm,yin2024online} & map, planner, often an LLM & detector & metric map & none & mostly categories\\
Learned ObjectNav policies~\cite{deitke2022procthor,chaplot2020object} & the policy & learned & learned or map & task-specific & categories\\
\midrule
\ours (this thesis) & metric controller & one small VLM & metric hypotheses, map & none & yes\\
\bottomrule
\end{tabular}}
\end{table}

\section{Semantic memory for embodied agents}
\label{sec:lit-memory}
Embodied agents attach semantics to geometry in many forms: open-vocabulary 3D scene graphs~\cite{gu2023conceptgraphs},
online semantic maps for changing scenes~\cite{jiang2025dualmap}, language fields and semantic Gaussian
splatting~\cite{katragadda2025online,lee2026embodiedsplat}, and long-horizon spatio-temporal memories for exploration,
question answering, and navigation~\cite{yang2024scene,anwar2024remembr,ginting2025enter}. These representations
are rich and general. They are built with dedicated encoders or large models, and they store far more than a
navigation agent needs to reach one object: a full scene graph or language field must be constructed and queried
at every step, which is costly on a small GPU and adds failure modes of its own. A navigation agent needs to know
where the target might be, how much evidence supports each candidate, and where it can move.

The memory of \ours also has a classical ancestor outside semantic mapping. Multiple-hypothesis
tracking keeps several candidate explanations of noisy, possibly false detections, associates each new detection
with one of them, and lets the evidence decide~\cite{reid1979algorithm}; probabilistic robotics applies the same
idea to localization and mapping~\cite{thrun2005probabilistic}. Active visual object search plans where a robot should
look for an object, using probabilistic models of where objects tend to be and how reliable its detector
is~\cite{aydemir2011search,aydemir2013active}. \ours borrows the hypothesis bookkeeping of tracking and the
detector-aware search of active object search, and replaces the trained detector of both with the yes/no and box
answers of a generative VLM, whose errors (confident boxes on look-alikes, repeated answers from one view) shape how
the evidence is counted.

Closest to the memory of \ours are two object-navigation systems. ApexNav fuses target-centric detections over time
to suppress false positives~\cite{zhang2025apexnav}, and AECNav treats expected but absent detections as negative
evidence~\cite{liu2026aecnav}. \ours shares both ideas and applies them at the minimal end of the spectrum: its memory
is a short list of metric target hypotheses whose evidence grows only with distinct viewpoints, a choice that
matters for a generative sensor whose repeated answers from the same view are not independent.

\section{Verification, uncertainty, and search}
\label{sec:lit-verify}
Spatially explicit prompting, such as marking candidate regions or actions in the image, improves the spatial
decisions of VLMs~\cite{yang2023mark,nasiriany2024pivot,goetting2024navigation}. Test-time sampling and verification,
adaptive allocation of test-time compute, and uncertainty-conditioned execution improve vision-language-action
models and embodied planners~\cite{kwok2025robomonkey,dao2026direct,choi2026scale}, and calibrated uncertainty lets
robots ask for help or explore until confident~\cite{ren2023robots,liang2024introspective,ren2024explore}.
These methods share one insight with \ours: a model's answer is more useful when the agent knows how much to
trust it. They differ in where the trust comes from. Sampling and verification spend more model calls on the same
decision; calibrated uncertainty asks the model how sure it is. \ours obtains trust from geometry: an answer gains
weight when the same location is confirmed from a different viewpoint, and a candidate is ruled out once the agent
reaches it without ending the task. It spends an extra query, a magnified crop, only where the evidence is weak.

\section{Geometry from images}
\label{sec:lit-geometry}
Learned reconstruction and depth networks recover metric geometry from RGB images: feed-forward reconstruction and
dense SLAM with learned priors~\cite{wang2025vggt,murai2024mast,maggio2025vggt}, zero-shot stereo~\cite{wen2025foundationstereo},
and monocular metric depth~\cite{yang2024depthv2,lin2025depth}. Sensor choices can change the conclusions of embodied
evaluations~\cite{corlito2026sensor,wang2025rethinking}, so an agent that claims to work from RGB must be evaluated
without privileged depth or pose. Monocular visual odometry fixes its unknown scale from the ground plane and the
known camera height~\cite{song2014robust}; \ours applies the same constraint to every frame of a monocular metric
depth network. Concurrently with this work, Li~\etal{} run Qwen3-VL-4B zero-shot and accumulate its answers as
per-direction scores in a bird's-eye field for vision-language navigation~\cite{li2026structure}; \ours instead lifts
boxes to metric target hypotheses for object navigation.

\section{Summary and research gap}
\label{sec:lit-gap}
Two lines of work meet in this thesis. MLLM agents follow open instructions but keep the model in charge of the plan,
and their gains depend on large models or training. Modular object-navigation systems keep a metric map and a
planner but rely on dedicated detectors, fixed vocabularies, and often large language models for exploration. No
existing agent lets a small generative VLM drive open-instruction navigation, and no study has isolated why such a
model fails when it is used as the agent. The thesis addresses both gaps: it diagnoses the failure
(\cref{chap:diagnosis}) and builds an agent in which the small VLM is only a semantic sensor (\cref{chap:method}).
